\section{Conclusion}
\label{sec:conclusion}

We have determined the complete automorphism group of the native
sixty-vertex quartic graph
\[
G60\cong \mathrm{AT4val}[60,6].
\]

The argument proceeds through two independent constructions.

First, the canonical fixed-point-free involution
\[
a\in\operatorname{Aut}(G60)
\]
defines a two-fold quotient
\[
G60\longrightarrow G30.
\]
The quotient automorphism group has order
\[
240,
\]
and every quotient automorphism has exactly two native lifts. This
constructs a subgroup of
\[
\operatorname{Aut}(G60)
\]
of order
\[
480.
\]

Second, an exhaustive enumeration performed directly on the exact
sixty-vertex graph returns exactly
\[
480
\]
automorphisms. The native enumeration and the lifted construction agree
as explicit permutation sets. Therefore
\[
\operatorname{Aut}(G60)
\]
is exactly the lifted group.

The independent enumeration also proves that every native automorphism
preserves the canonical $a$-fiber system and centralizes the deck
involution. Thus the quotient is not merely a convenient presentation.
It is an invariant structure respected by the full automorphism group.

The internal group invariants are
\[
Z(\operatorname{Aut}(G60))\cong C_2,
\]
\[
\operatorname{Aut}(G60)'
\cong
A_5\times C_2,
\]
\[
\operatorname{Aut}(G60)''
\cong
A_5,
\]
and
\[
\operatorname{Aut}(G60)^{\mathrm{ab}}
\cong
C_2\times C_2.
\]

The full element-order profile and the three index-two subgroups identify
the correct abstract model as the fiber product
\[
\operatorname{Aut}(G60)
\cong
S_5\times_{C_2}D_8,
\]
where the sign character of $S_5$ is matched with the quotient character
of $D_8$ having Klein four kernel.

This identification is not based only on matching invariants. We
construct an explicit bijection between the $480$ native automorphisms
and the $480$ elements of the fiber-product model. The map preserves all
\[
480^2=230400
\]
ordered products.

The final result may therefore be stated in three equivalent forms.

\begin{theorem}
For
\[
G60\cong \mathrm{AT4val}[60,6],
\]
the following hold.

\begin{enumerate}
\item The full automorphism group has order
\[
\left|\operatorname{Aut}(G60)\right|=480.
\]

\item Every automorphism of $G60$ is a lift of an automorphism of the
canonical thirty-vertex quotient.

\item The full group is explicitly isomorphic to
\[
S_5\times_{C_2}D_8,
\]
using the Klein-four-kernel quotient character of $D_8$.
\end{enumerate}
\end{theorem}

The classification problem is therefore closed.

What remains is interpretive rather than identificatory. The dihedral
coordinate should be understood geometrically in the native action: the
central deck involution, the order-four directions, the reflections, and
the parity coupling to the $S_5$ coordinate all admit further geometric
explanation.

Those questions concern why the symmetry group has this architecture.
They no longer concern which group it is.
